\subsection{Аппаратная часть: 4 + 25 светодиодов}

На плате Пионера четыре RGB-светодиода --- по одному у каждого луча. Модуль светодиодной матрицы добавляет ещё 25. Программа видит их как одну линейку: индексы \verb|0..3| --- светодиоды платы, \verb|4..28| --- модуль. Всего 29, поэтому в курсе линейка всегда создаётся как \verb|Ledbar.new(29)|: на дроне без модуля лишние индексы просто ни на что не влияют, а в симуляторе отображаются все 29.

\begin{center}
\begin{tikzpicture}[scale=1.1]
    \draw[fill=gray!10, thick, rounded corners] (-1.5,-1.5) rectangle (1.5,1.5);
    \node at (0,0.35) {\textbf{Pioneer}};
    \node[font=\footnotesize] at (0,-0.15) {(вид сверху)};
    \draw[->, ultra thick, GeoskanBlue] (0, 1.6) -- (0, 2.6) node[above, font=\footnotesize\bfseries] {Вперёд, ось X};
    \foreach \x/\y in {1.5/1.5, 1.5/-1.5, -1.5/1.5, -1.5/-1.5} { \draw[thick, double] (\x,\y) -- (1.6*\x,1.6*\y); }
    \foreach \x/\y in {2.4/2.4, 2.4/-2.4, -2.4/2.4, -2.4/-2.4} { \draw[fill=black!15] (\x,\y) circle (0.55); }
    \node[circle, fill=red!70, inner sep=3pt, label=left:{\ttfamily 0}] at (-1.1, 1.1) {};
    \node[circle, fill=green!70!black, inner sep=3pt, label=right:{\ttfamily 1}] at (1.1, 1.1) {};
    \node[circle, fill=blue!70, inner sep=3pt, label=right:{\ttfamily 2}] at (1.1, -1.1) {};
    \node[circle, fill=yellow!80!black, inner sep=3pt, label=left:{\ttfamily 3}] at (-1.1, -1.1) {};
    \node[font=\footnotesize, align=center] at (0,-3.4) {Индексы \texttt{0..3} --- плата; \texttt{4..28} --- модуль матрицы.\\ Какой индекс какому углу соответствует, проверяет задание ниже.};
\end{tikzpicture}
\end{center}

\subsubsection*{Модель RGB}

Цвет --- три числа от \verb|0.0| до \verb|1.0|: красный, зелёный, синий. Смешивая каналы, получаем любой оттенок:

\begin{center}
\small
\begin{tabular}{@{}llll@{}}
\toprule
\verb|{1, 0, 0}| красный & \verb|{0, 1, 0}| зелёный & \verb|{0, 0, 1}| синий & \verb|{1, 1, 0}| жёлтый \\
\verb|{1, 0, 1}| фиолетовый & \verb|{0, 1, 1}| бирюзовый & \verb|{1, 1, 1}| белый & \verb|{0, 0, 0}| выключен \\
\bottomrule
\end{tabular}
\end{center}

Значения меньше единицы --- тусклее: \verb|{0.2, 0, 0}| --- тёмно-красный. В Python светодиодами управляет метод \verb|led_control| библиотеки \texttt{pioneer\_sdk}, и компоненты там --- целые числа \verb|0..255| (глава~\ref{sec:python-sdk}); в Lua --- только \verb|0..1|.

\subsection{Модуль Ledbar}

\begin{codebox}[Две функции, которых достаточно]
\begin{itemize}
  \item \verb|local leds = Ledbar.new(29)| --- создать объект линейки из 29 светодиодов. Вызывается один раз в начале программы.
  \item \verb|leds:set(i, r, g, b)| --- задать цвет светодиода \verb|i| (от \verb|0| до \verb|28|). Обратите внимание на двоеточие: \verb|set| --- метод объекта \verb|leds|.
\end{itemize}
\end{codebox}

\begin{codefile}{lua/05_leds/leds_basic_green.lua}
\end{codefile}

\paragraph{Что происходит.} Функция \verb|changeColor| принимает таблицу \verb|{r, g, b}|, а \verb|unpack| раскладывает её в три аргумента \verb|leds:set|. Цикл \verb|for i = 0, ledNumber - 1| обходит все 29 индексов. Один вызов \verb|changeColor({0, 1, 0})| --- и вся линейка зелёная.

\begin{warnbox}[Индекс вне диапазона не вызывает ошибку]
\verb|leds:set(29, ...)| на линейке из 29 элементов молча игнорируется: светодиод не загорится, а программа продолжит работать. Если какой-то светодиод «не светится», первым делом проверьте индекс.
\end{warnbox}

\subsubsection*{Исследование: какой индекс где}

\begin{codefile}{lua/05_leds/leds_addressing_research.lua}
\end{codefile}

\begin{taskbox}[Карта светодиодов]
Запустите программу: загорится светодиод с индексом \verb|0|. Посмотрите на дрон (в симуляторе --- камерой \ui{Дрон}): какой это угол? Замените \verb|0| на \verb|1|, \verb|2|, \verb|3| и заполните таблицу. Если есть модуль матрицы --- продолжите до \verb|28|.
\end{taskbox}

\begin{center}
\renewcommand{\arraystretch}{1.4}
\begin{tabular}{|c|p{7cm}|}
\hline
\rowcolor{GeoskanBlue!20} \textbf{Индекс} & \textbf{Расположение} \\
\hline
0 & \dotfill \\ \hline
1 & \dotfill \\ \hline
2 & \dotfill \\ \hline
3 & \dotfill \\ \hline
4 \dots 28 & \dotfill \\ \hline
\end{tabular}
\end{center}

Это упражнение учит главному принципу инженера: если документации недостаточно --- проведи эксперимент.
